\title{The Era of 1-bit LLMs: All Large Language Models are in 1.58 Bits}

\begin{abstract}
Recent advancements, exemplified by BitNet~\cite{bitnet}, are ushering in a new age of 1-bit Large Language Models (LLMs). In this paper, we present a 1-bit LLM variant, \textbf{\text{BitNet b1.58}}, where every individual parameter (or weight) in the LLM takes on a ternary value \{-1, 0, 1\}. This model matches the full-precision (i.e., FP16 or BF16) Transformer LLM of identical size and training data in terms of both perplexity and downstream task performance, while offering considerable improvements in latency, memory usage, throughput, and energy efficiency. More importantly, the 1.58-bit LLM introduces a novel scaling law and training methodology for developing next-generation LLMs that are simultaneously high-performing and cost-efficient. Additionally, it fosters a new computational paradigm and paves the way for the creation of specialized hardware tailored to 1-bit LLMs.
\end{abstract}

\section{The Era of 1-bit LLMs}

The domain of AI has experienced rapid expansion in the size and capabilities of Large Language Models (LLMs) in recent years. These models have showcased outstanding effectiveness across various natural language processing tasks; however, their growing scale has introduced challenges for deployment and sparked concerns over their environmental and financial footprint due to substantial energy demands.

A strategy to mitigate these issues involves post-training quantization, which generates low-bit models optimized for inference~\cite{smoothquant, gptq, quip, quip_sharp}. This approach lowers the precision of weights and activations, significantly cutting down the memory and computational overhead of LLMs. The progression has moved from 16-bit representations towards even lower-bit formats, such as 4-bit versions~\cite{gptq, awq}. Nonetheless, post-training quantization remains suboptimal, despite its widespread adoption in industrial LLM deployments.

Recent research on 1-bit model architectures, such as BitNet~\cite{bitnet}, offers a promising avenue for lowering the expenses associated with LLMs while preserving their performance levels. Standard LLMs operate with 16-bit floating-point values (i.e., FP16 or BF16), and the majority of computations in these models consist of matrix multiplications. Hence, the primary computational cost arises from floating-point addition and multiplication operations. In contrast, BitNet's matrix multiplication involves only integer addition, resulting in substantial energy savings for LLMs. Since power is a fundamental constraint on computational performance in many chips, these energy savings can translate into faster processing speeds.

Beyond computation, transferring model parameters from DRAM to the memory of an on-chip accelerator (such as SRAM) can be costly during inference. While attempts have been made to expand SRAM capacity to boost throughput, this approach significantly increases costs compared to DRAM. Relative to full-precision models, 1-bit LLMs have a considerably smaller memory footprint both in terms of capacity and bandwidth requirements. This reduction can greatly decrease the cost and latency of loading weights from DRAM, enabling faster and more efficient inference.

In this paper, we present a notable 1-bit LLM variant named \textbf{\text{BitNet b1.58}}, where each parameter is ternary, assuming values in \{-1, 0, 1\}. We introduce an additional zero value to the original 1-bit BitNet, resulting in a representation size of 1.58 bits in the binary system. \text{BitNet b1.58} maintains all the advantages of the original 1-bit BitNet, including its novel computational paradigm that nearly eliminates multiplication operations in matrix multiplications and is amenable to extensive optimization. Moreover, it shares the same energy consumption profile as the original 1-bit BitNet and significantly improves memory usage efficiency, throughput, and latency compared to FP16 LLM baselines. Additionally, \text{BitNet b1.58} offers two further benefits. First, its modeling power is enhanced due to explicit feature filtering enabled by incorporating zero weights, which can substantially boost the performance of 1-bit LLMs. Second, our experimental results demonstrate that \text{BitNet b1.58} can match full-precision (i.e., FP16) baselines in both perplexity and downstream task performance starting at a model size of 3B, using identical configurations (e.g., model size, training tokens, etc.).

\section{BitNet b1.58}

\text{BitNet b1.58} builds upon the BitNet framework, which is a Transformer architecture that substitutes \emph{nn.Linear} with \emph{BitLinear}. It is trained from the ground up, employing 1.58-bit precision for weights and 8-bit precision for activations. Compared to the original BitNet, several adjustments have been made, which we outline below.

\paragraph{Quantization Function.} To restrict the weights to values of -1, 0, or +1, we utilize an \emph{absmean} quantization approach. This method scales the weight matrix by its mean absolute value and then rounds each entry to the closest integer within \{-1, 0, +1\}:

\begin{equation}
    \widetilde{W} = \mathrm{RoundClip}(\frac{W}{\gamma + \epsilon}, -1, 1),
\end{equation}

\begin{equation}
    \mathrm{RoundClip}(x, a, b) = \max(a, \min(b, \mathrm{round}(x))),
\end{equation}

\begin{equation}
    \gamma = \frac{1}{nm}\sum_{ij} |W_{ij}|.
\end{equation}

For activations, the quantization procedure follows the same design as in BitNet, except that we omit scaling the activations before nonlinearities to the interval $[0, Q_b]$. Instead, activations are scaled per token to lie within $[-Q_b, Q_b]$, thereby eliminating the need for zero-point quantization. This simplifies implementation and system-level optimization without notably impacting performance in our experiments.

\paragraph{LLaMA-inspired Components.}

The LLaMA architecture~\cite{llama, llama2} has become the standard backbone for open-source LLMs. To align with the open-source community, \text{BitNet b1.58} incorporates LLaMA-like elements. Specifically, it adopts RMSNorm~\cite{rmsnorm}, SwiGLU~\cite{swiglu}, rotary embeddings~\cite{rope}, and removes all bias terms. This design facilitates easy integration of \text{BitNet b1.58} into widely used open-source platforms such as Huggingface, vLLM~\cite{vllm}, and llama.cpp\footnote{\url{https://github.com/ggerganov/llama.cpp}} with minimal modification.

\section{Results}

We conducted a comparison between \text{BitNet b1.58} and our reproduced FP16 \text{LLaMA LLM} across different model sizes. To maintain fairness, both models were pre-trained on the RedPajama dataset~\cite{redpajama} using 100 billion tokens. Their zero-shot performance was assessed on multiple language tasks, including ARC-Easy~\cite{arc}, ARC-Challenge~\cite{arc}, Hellaswag~\cite{hellaswag}, Winogrande~\cite{winoGrande}, PIQA~\cite{piqa}, OpenbookQA~\cite{openbookqa}, and BoolQ~\cite{boolq}. Additionally, we reported validation perplexity results on the WikiText2~\cite{wikitext2} and C4~\cite{c4} datasets.

We evaluated the runtime GPU memory usage and latency for both \text{LLaMA LLM} and \text{BitNet b1.58}. These metrics were obtained using the FasterTransformer\footnote{\url{https://github.com/NVIDIA/FasterTransformer}} framework, which is optimized for minimizing LLM inference latency on GPU platforms. The 2-bit kernel from Ladder~\cite{ladder} was also incorporated for \text{BitNet b1.58}. We reported the inference time per output token, as this represents the primary computational cost during inference.

Table~\ref{tab:ppl} provides a summary of perplexity and resource costs for \text{BitNet b1.58} and \text{LLaMA LLM}. It indicates that \text{BitNet b1.58} attains parity with full precision \text{LLaMA LLM} at a 3B model scale in terms of perplexity, while achieving 2.71 times faster speed and utilizing 3.55 times less GPU memory. Notably, the 3.9B variant of \text{BitNet b1.58} runs 2.4 times faster, uses 3.32 times less memory, and significantly outperforms the 3B \text{LLaMA LLM}.

Detailed zero-shot accuracy results on downstream tasks are presented in Table~\ref{tab:zero-shot}. We employed the \emph{lm-evaluation-harness}\footnote{\url{https://github.com/EleutherAI/lm-evaluation-harness}} evaluation pipeline. The findings reveal that the performance gap between \text{BitNet b1.58} and \text{LLaMA LLM} diminishes as model size grows. Crucially, \text{BitNet b1.58} matches the full precision baseline starting at 3B. Echoing the perplexity trends, the final task evaluations show that the 3.9B \text{BitNet b1.58} surpasses the 3B \text{LLaMA LLM} while demanding less memory and lower latency. This confirms that \text{BitNet b1.58} constitutes a Pareto improvement over current leading LLM models.

\paragraph{Memory and Latency}

We additionally increased the model sizes to 7B, 13B, and 70B to assess the associated costs. Figure~\ref{fig:latency-memory-energy} displays the latency and memory trends, demonstrating that speed-up improves as the model size grows. Notably, \text{BitNet b1.58} 70B achieves a speed 4.1 times greater than the \text{LLaMA LLM} baseline. This improvement arises because the time required for \emph{nn.Linear} scales with the model size. Memory usage shows a similar pattern, since the embedding remains in full precision and its memory share diminishes for larger models. Both latency and memory were evaluated using a 2-bit kernel, indicating potential for further cost reduction through optimization.

\paragraph{Energy}

We also estimated the energy consumption of arithmetic operations for both \text{BitNet b1.58} and \text{LLaMA LLM}. Our analysis primarily concentrates on matrix multiplication calculations, as they dominate the cost in LLMs. Figure~\ref{fig:energy} depicts the breakdown of energy expenses. Most of \text{BitNet b1.58}'s energy usage comes from INT8 addition operations, whereas \text{LLaMA LLM} involves both FP16 addition and FP16 multiplication. Based on the energy model in~\cite{energycost, pokebnn}, \text{BitNet b1.58} reduces arithmetic operation energy consumption for matrix multiplication by a factor of 71.4 on 7nm chips. We also present the overall energy cost for processing models with 512 tokens. Our findings indicate that as model size increases, \text{BitNet b1.58} becomes progressively more energy efficient relative to the FP16 \text{LLaMA LLM} baseline. This efficiency gain results from the increasing dominance of \emph{nn.Linear} operations with larger models, while the energy cost from other components decreases proportionally.

\paragraph{Throughput}

We evaluated the throughput of \text{BitNet b1.58} and the \text{LLaMA LLM} with 70 billion parameters on two 80GB A100 GPUs, utilizing pipeline parallelism~\cite{gpipe} to enable running \text{LLaMA LLM} 70B on these devices. The batch size was progressively increased until reaching the GPU memory capacity, with a fixed sequence length of 512. As shown in Table~\ref{tab:throughput}, \text{BitNet b1.58} 70B supports up to 11 times larger batch sizes than \text{LLaMA LLM}, resulting in an 8.9-fold increase in throughput.

\textbf{\text{BitNet b1.58} establishes a new scaling paradigm concerning model performance and inference cost}. For reference, based on the findings in Figures~\ref{fig:latency-memory-energy} and \ref{fig:energy}, the following equivalences hold between models of different sizes in 1.58-bit versus 16-bit:

\begin{itemize}
    \item A 13B \text{BitNet b1.58} is more efficient with respect to latency, memory usage, and energy consumption compared to a 3B FP16 LLM.
    \item A 30B \text{BitNet b1.58} surpasses a 7B FP16 LLM in latency, memory footprint, and energy efficiency.
    \item A 70B \text{BitNet b1.58} outperforms a 13B FP16 LLM in terms of latency, memory demands, and energy consumption.
\end{itemize}

\paragraph{Training with 2T Tokens}

The total number of training tokens is a vital variable for large language models. To assess the token scalability of \text{BitNet b1.58}, we trained a \text{BitNet b1.58} model on 2 trillion tokens using the data preparation approach from StableLM-3B~\cite{StableLM-3B-4E1T}, which represents the leading open-source 3B model currently available. Both models were benchmarked on a suite including Winogrande~\cite{winoGrande}, PIQA~\cite{piqa}, SciQ~\cite{sciq}, LAMBADA~\cite{lambada}, and ARC-easy~\cite{arc}. Zero-shot accuracy scores are presented in Table~\ref{tab:2t}. For tasks evaluated via accuracy and normalized accuracy, we report their average. The performance metrics for StableLM 3B at 2T tokens were obtained directly from its official technical report. Our results demonstrate that \text{BitNet b1.58} attains superior outcomes across all evaluated tasks, suggesting that 1.58-bit LLMs possess robust generalization abilities.

\section{Discussion and Future Work}

\textbf{1-bit Mixture-of-Experts (MoE) LLMs}

Mixture-of-Experts (MoE) models have shown to be a cost-efficient strategy for large language models (LLMs). Although they substantially decrease computational FLOPs, the high memory usage and the overhead from inter-chip communication constrain their deployment and practical use. These issues can be mitigated by 1.58-bit LLMs. Firstly, the diminished memory requirement lowers the number of devices necessary to deploy MoE models. Additionally, it greatly reduces the cost of transferring activations across networks. Ultimately, if the entire model fits onto a single chip, this overhead would be eliminated.

\textbf{Native Support of Long Sequence in LLMs}

With the rise of LLMs, handling long sequences has become an essential requirement. A significant obstacle in long sequence inference is the memory demand caused by KV caches. \text{BitNet b1.58} marks an important advancement toward native long sequence support by cutting activations from 16 bits down to 8 bits, enabling the context length to be doubled with the same resources. This can be further compressed without loss to 4 bits or even lower for 1.58-bit LLMs, which we consider as future work.

\textbf{LLMs on Edge and Mobile}

Implementing 1.58-bit LLMs could dramatically enhance the performance of language models on edge and mobile platforms. Such devices often face constraints in memory and computational power, limiting the scale and efficiency of LLMs. The lowered memory and energy demands of 1.58-bit LLMs make it feasible to deploy these models on edge and mobile devices, unlocking numerous applications that were previously infeasible. This advancement can significantly boost the capabilities of edge and mobile devices and enable novel, exciting uses of LLMs. Furthermore, 1.58-bit LLMs are more compatible with CPU-based devices, which are predominantly used in edge and mobile environments. Thus, \text{BitNet b1.58} can be run efficiently on such devices, further enhancing their performance and functionalities.

\textbf{New Hardware for 1-bit LLMs}

Recent developments such as~Groq\footnote{\url{https://groq.com/}} have shown encouraging results and significant promise in creating specialized hardware (e.g., LPUs) tailored for LLMs. Taking this further, we anticipate and advocate for designing new hardware and systems specifically optimized for 1-bit LLMs, taking advantage of the new computational paradigm introduced by BitNet~\cite{bitnet}.

\end{document}